% Einheit 3 von 4 – Flächen- und Volumeneinheiten (bank/einheiten/e3.jsonl, zusammenbau v0.1)
%% TODO Zweigzeile Teil 1: was hier gelernt wird (Satz in Schülersprache aus dem Einheitstitel) – Einheit 3
\einheitenkopf[e3]{Einheit 3 von 4 · Flächen- und Volumeneinheiten}
\zweigzeile{ab Kl. 5 · P10}
\setcounter{aufgabe}{28}

%% TODO \verfahren{…}: Verfahrensüberschrift in Sachsprache fehlt in der Bank (2.3 b)
%% TODO Titel: Ich-kann-Satz fehlt in der Bank (Platzhalter: Kettenname) – Nr. 29
%% TODO Anweisung: ein Satz über der ersten Teilaufgabe fehlt in der Bank – Nr. 29
\begin{aufgabe}{Flächen- und Volumeneinheiten umrechnen}
%% TODO \rechenplatz{n}: Schrittzahl des Grundfalls fehlt in der Bank (nur bei fester Schreibform, 2.3 b)
\begin{teile}
\teil Quadratmeter und Quadratdezimeter – welche Umrechnungszahl? Kreuze an.\\ \kreuz{zehn}\\ \kreuz{hundert}\\ \kreuz{tausend}
\teil $3$ m² – wie viel dm²? \leerfeld dm²
\teil $4$ m² – wie viel cm²? \leerfeld[cm²]
\teil $7$ a – wie viel m²? \leerfeld[m²]
\teil Ordne, mit der kleinsten beginnend: m³, mm³, dm³, cm³ \leerfeld ; \leerfeld ; \leerfeld ; \leerfeld
\teil $5$ dm³ – wie viel cm³? \leerfeld[cm³]
\teil $6$ l – wie viel dm³? \leerfeld dm³
\teil $1{,}8$ l – wie viel ml? \leerfeld[ml]
\teil Was bedeutet „Kilo“ in Kilowatt? Kreuze an.\\ \kreuz{tausend}\\ \kreuz{hundert}\\ \kreuz{ein Tausendstel}
\teil Eine Flasche enthält $1{,}25$ l Saft. Wie viele Gläser zu $200$ ml kann man damit ganz füllen? \hfill (P10 2022 OS) \\ \leerfeld Gläser
\teil Eine Pumpe fördert $12\,500$ l pro Stunde. Wie lange braucht sie, um einen Teich mit $75$ m³ Wasser zu füllen? \hfill (P10 2017 OS) \\ \leerfeld[h]
\end{teile}
\end{aufgabe}

%% TODO \verfahren{…}: Verfahrensüberschrift in Sachsprache fehlt in der Bank (2.3 b)
%% TODO Titel: Ich-kann-Satz fehlt in der Bank (Platzhalter: Kettenname) – Nr. 30
%% TODO Anweisung: ein Satz über der ersten Teilaufgabe fehlt in der Bank – Nr. 30
\begin{aufgabe}{Flächen- und Volumeneinheiten im Koordinatensystem, Sek II}
%% TODO \rechenplatz{n}: Schrittzahl des Grundfalls fehlt in der Bank (nur bei fester Schreibform, 2.3 b)
\begin{teile}
\teil Eine Längeneinheit entspricht $15$ cm; gesucht ist der Flächeninhalt eines Schildes. Was wird umgerechnet? Kreuze an.\\ \kreuz{Länge}\\ \kreuz{Fläche}\\ \kreuz{Volumen}
\teil Eine Fläche ist $6$ Flächeneinheiten groß; eine Längeneinheit entspricht $5$ cm. Wie viel cm²? \leerfeld[cm²]
\teil Die Vorderseite einer Gartenhütte ist ein Rechteck, $3$ Längeneinheiten breit und $2$ hoch, mit einem Dachdreieck darauf (Grundseite $3$, Höhe $1{,}5$). Eine Längeneinheit entspricht einem Meter, die Hütte ist $4$ m tief. Querschnitt und Volumen? \leerfeld[m²] , \leerfeld[m³]
\teil Ein Aufsteller hat die Form der Fläche zwischen $g(x) = -0{,}25x^2 + 2$ (oben) und $f(x) = 0{,}25x^4 - x^2 + 1$ (unten); die Schnittstellen liegen bei $x = -2$ und $x = 2$, eine Längeneinheit entspricht $30$ cm. Weise nach, dass die Fläche $4{,}8$ Flächeneinheiten groß ist. Der Aufsteller wird auf beiden Seiten beklebt. Wie viel Folie braucht man in cm² und in m²? \hfill (FHR 2022) \\ \leerfeld[cm²] , \leerfeld[m²]
\end{teile}
\end{aufgabe}

%% TODO Titel: Ich-kann-Satz fehlt in der Bank (Platzhalter: Kettenname) – Nr. 31
%% TODO Anweisung: ein Satz über der ersten Teilaufgabe fehlt in der Bank – Nr. 31
\begin{aufgabe}{Flächeneinheiten ordnen (mm² bis km², a und ha einordnen)}
\begin{teile}
\teil Ordne, mit der kleinsten beginnend: m², ha, cm², a, km² \leerfeld ; \leerfeld ; \leerfeld ; \leerfeld ; \leerfeld
\end{teile}
\end{aufgabe}

%% TODO Titel: Ich-kann-Satz fehlt in der Bank (Platzhalter: Kettenname) – Nr. 32
%% TODO Anweisung: ein Satz über der ersten Teilaufgabe fehlt in der Bank – Nr. 32
\begin{aufgabe}{Repräsentanten zuordnen}
\begin{teile}
\teil Ordne zu: Briefmarke, Kinderzimmer, Fußballfeld – $7\,000$ m²; $6$ cm²; $18$ m² Briefmarke \leerfeld , Zimmer \leerfeld , Fußballfeld \leerfeld
\end{teile}
\end{aufgabe}

%% TODO Titel: Ich-kann-Satz fehlt in der Bank (Platzhalter: Kettenname) – Nr. 33
%% TODO Anweisung: ein Satz über der ersten Teilaufgabe fehlt in der Bank – Nr. 33
\begin{aufgabe}{Flächen- und Volumeneinheiten umrechnen – Fehler finden, Begründen, Anwendung, Darstellungswechsel}
\begin{teile}
\teil Nils schreibt: $7$ dm² $= 70$ cm². Wo ist der Fehler?
\teil Eine Längeneinheit entspricht $25$ cm, eine Fläche ist $8$ Flächeneinheiten groß. Ole rechnet: $8 \cdot 25 = 200$ cm². Wo ist der Fehler?
\teil Begründe: Warum passen in einen Quadratmeter hundert Quadratdezimeter und nicht zehn?
\teil Eine Längeneinheit entspricht $3$ cm. Begründe, warum eine Flächeneinheit nicht $3$ cm², sondern $9$ cm² entspricht.
\teil Ein Aquarium ist innen $90$ cm lang, $40$ cm breit und $45$ cm hoch. Es wird bis $5$ cm unter den Rand gefüllt. Wie viele Liter Wasser braucht man? \leerfeld[l]
\teil Trage $3{,}0715$ m² in die Einheitentafel ein (zwei Stellen je Einheit) und gib die Fläche in cm² an.

\sachtabelle{ccc}{m² & dm² & cm²}{\leerzelle & \leerzelle & \leerzelle}
\leerfeld[cm²]
\end{teile}
\end{aufgabe}
